% Learning InitechBase, Chapter 7: Reports (AN-16).  Every panel replayed
% (examples/l7-*.dot); sources cmd/report.c (REPORT FORM, mint 004 layout),
% the report forms STAFFRPT.FRM and INVRPT.FRM of the sample diskette.
\chapter{Reports}
\Gates{test-pubs-examples, test-samir-report}
A report arranges the records of a database in columns, under a title and
headings, with totals.  The layout of a report is kept in a report form file,
so the same report can be produced again whenever the data changes.  In this
lesson you produce the staff salary report of Initech from the report form on
the sample diskette.

\begin{ChapterCovers}
\item Report forms
\item \cmd{REPORT FORM}
\item Groups and subtotals
\item Choosing records for a report
\item \kw{PLAIN}, \kw{SUMMARY} and \kw{HEADING}
\item Writing a report to a file
\end{ChapterCovers}

\ATEHead{Preparing for This Chapter}
Start InitechBase with a fresh copy of the sample files.

\ATEHead{The Report Form}
The report form \file{STAFFRPT.FRM} describes the staff salary report.  Its
title has two lines; it has four columns -- the last name, the first name, the
date of hire and the salary -- of which the salary is totalled; and it groups
the employees by department, with a subtotal for each.  A report form is used
with the database whose fields it names.

The groups follow one another as the records come, so the records must be in
the order of the group: open the staff database with the department index.

\ATEHead{Producing the Report}
\begin{Steps}
\item Type \type{USE STAFF INDEX STAFFDPT}.
\item Type \type{REPORT FORM STAFFRPT}.
\ATEResponse{InitechBase prints the report on the screen.}
\end{Steps}
\DotExample{examples/l7-report}
Each page begins with its number, the date and the title, then the column
headings.  Each group begins with its heading and ends with its subtotal; the
total of all the groups ends the report.  Robert Chen, whose record is marked
for deletion, is in the report: \cmd{REPORT FORM} includes marked records
unless SET DELETED is ON.

\ATEHead{Choosing the Records}
\kw{FOR} chooses the records of a report as it does for \cmd{LIST}.
\kw{PLAIN} leaves out the page number and the date.
\begin{Steps}[resume]
\item Type \type{REPORT FORM STAFFRPT FOR EXEMPT PLAIN}.
\end{Steps}
\DotExample{examples/l7-plain}

\ATEHead{Totals Only}
\kw{SUMMARY} prints the group headings and the totals without the lines of the
records:
\DotExample{examples/l7-summary}

\ATEHead{A Heading of Your Own}
\kw{HEADING} adds a line of your own after the page number.  Here the report of
the Accounting department is produced for the month, with the marked records
left out:
\begin{Steps}[resume]
\item Type \type{SET DELETED ON}.
\item Type \type{REPORT FORM STAFFRPT HEADING 'MARCH 1992' FOR DEPT = 'ACCT'}.
\end{Steps}
\DotExample{examples/l7-heading}

\ATEHead{Writing a Report to a File}
\kw{TO FILE} writes the report to a text file as well as displaying it.  You can
copy the file to a printer, or keep it, from InitechDOS.  The report form
\file{INVRPT.FRM} is the open invoice register of the accounts receivable
ledger:
\begin{Steps}[resume]
\item Type \type{USE INVOICE INDEX INVNO}.
\item Type \type{REPORT FORM INVRPT PLAIN TO FILE INVRPT}.
\item Type \type{? FILE('INVRPT.TXT')}.
\ATEResponse{\fn{FILE()} is true: the file \file{INVRPT.TXT} is on the disk.}
\end{Steps}
\DotExample{examples/l7-tofile}

\ATEHead{Where Do You Go From Here?}
You can produce a report from a report form, for all the records or some of
them.  In Chapter~8 you work with two databases at once.

\ATEHead{Summary}
\begin{enumerate}
\item \cmd{REPORT FORM} \xp{form} prints the report described by a report form
  file.
\item For a grouped report, open the database with an index in the order of the
  group.
\item \kw{FOR} and \kw{WHILE} choose the records; SET DELETED ON leaves out the
  marked records.
\item \kw{PLAIN}, \kw{SUMMARY} and \kw{HEADING} change the layout.
\item \kw{TO FILE} \xp{file} writes the report to a text file.
\end{enumerate}
